\begin{nestedsection}{CoDeR Model for Continuous Deductive Reasoning}{model}
	CoDeR is an abstract model for continuous deductive reasoning over streamed RDF data that expresses the semantics of Continuous Datalog, thus supporting continuous reasoning expressed in both RIF-Core and OWl 2 RL.
	It is composed of a stream-based data model for expressing the streams $S$, $P$ and $N$ of each ${CDP = \{A,S,R\} \vDash \{A',P,N\}}$, and a minimal algebra of stream-to-stream operators for expressing the persistent set of axiomatic rules $A$.
	It casts the problem of applying those rules to streamed RDF data as that of iterative pattern matching, in a manner inspired by the Rete pattern matching algorithm \citep{forgy79}, though it is meant to support any such algorithm rather than prescribing one itself.

	% Stream-based Data Model
	\begin{nestedsection}{Data Model}{model: data}
		The data model of CoDeR is composed of two classes of streams: \emph{input}/\emph{output} streams and \emph{inter-operator} streams.

		In the first case, input/output streams provide the semantics of the streams $S$, $P$ and $N$ for any ${CDP}$ expressed by a CoDeR-based system;
		they are sequences of instances of RDF triples, where an instance of an RDF triple is an RDF triple annotated with the time interval for which it is valid \citep{SemanticStreamingManagement,sparkwave} according to the ${CDP}$ expressed: ${\langle (s,p,o),i_{e},i_{n} \rangle}$.
		This interval, for any entailed fact instance, is derived from the valid time intervals of the fact instances that justify it, as shown in \reffig{intersected-intervals}, being the graphical representation of valid times for transitive data entailed by the state of the ${CDP = \{\{c\,\text{:-}\,a\,,\,b\,.\},S,3\}}$ at instant ${i = 2}$:
		\begin{figure}[t]
			\centering
			\includegraphics[width=0.45\textwidth]{intersected-intervals}
			\caption{Derivation of valid-time for entailed fact instance from those of its justifying fact instances.}
			\labelfig{intersected-intervals}
		\end{figure}
		\begin{multline*}
			CDP_{2} = \{c\,\text{:-}\,a\,,\,b\,.\} \cup \{ \langle a,1,4 \rangle . \langle b,2,5 \rangle \} \vDash \\
			\{ c\,\text{:-}\,a\,,\,b\,.\,\langle a,1,4 \rangle . \langle b,2,5 \rangle . \langle c,2,4 \rangle \}
		\end{multline*}
		It is in the base case of this derivation that the semantic window range $R$ of a ${CDP}$ is expressed:
		each fact instance arriving on an input stream has a valid interval from the time of its arrival at the system for a duration equal to the range $R$.

		In the second case, inter-operator streams are produced and consumed by the stream-to-stream operators of CoDeR's minimal algebra.
		These streams are sequences of the intermediate results of processing between one step in the production of entailments and the next.
		As the production of entailments in CoDeR is cast as a pattern matching problem, all intermediate results of this process will take the form of instances of RDF \emph{graphs};
		for a given inter-operator stream, each such RDF graph will match the same graph pattern, that being a sub-pattern of the body of one of the axiomatic rules $A$.
		As with instances of RDF triples, instances of RDF graphs are RDF graphs that are annotated with the time interval for which they are ``valid'': ${\langle \{(s,p,o),\dots,(s',p',o')\},i_{e},i_{n} \rangle}$.
		As these instances of RDF graphs may be considered to be ``entailed'' by the presence of valid instances of their constituent RDF triples within the system, the same semantics of validity may be applied to these intermediate matches as to the true entailments of the model, and their valid-time intervals may derived from their justifying instances in the same manner.
		It follows that inter-operator streams are, functionally, a general case of input/output streams, as the latter are simply sequences of instances of singleton RDF graphs matching the pattern ${\{(?s,?p,?o)\}}$.

		Given that the ordering of instances in the output stream of a CoDeR-based system by the value of $i_{e}$ yields the sequence of instances in the semantic stream $P$ entailed by a ${CDP}$, and that their ordering by the value of $i_{n}$ yields the sequence of instances in the semantic stream $N$, this gives the output streams of CoDeR a dual semantics with respect to Continuous Datalog.
		Furthermore, this dual semantics applies to all the streams of CoDeR, each expressing both the stream of entailment instances \emph{and} instance negations produced by a specific pipeline of CoDeR operators applied to some set of continuous RDF data sources.
	\end{nestedsection}

	% Minimal Algebra


%	\subimport{conclusions/}{future-work.tex}
\end{nestedsection}